\chapter{Knowledge Base and Data Preparation}
\label{chap:knowledge-base}\index{knowledge base}

The assistant can only say what its knowledge base contains, and it can only cite what the
knowledge base lets it locate. Preparing the data is therefore not a technical detail: it decides
which questions can be answered, how precisely, and with which source. The preparation is done
offline, before the assistant is started, and follows the order of this chapter: the documents
are selected (Section~\ref{sec:knowledge-base}), made synthetic (Section~\ref{sec:synthetic-data}),
converted from \ac{pdf} into structured text (Section~\ref{sec:parsing}), cleaned of what is not
text (Section~\ref{sec:non-text}) and anonymized (Section~\ref{sec:anonymization}).

\section{Sources and Types of Data}
\label{sec:knowledge-base}\index{FAQ}\index{forms}\index{DIP}\index{CGA}\index{support pages}

The public part of the knowledge base comes from public insurance documentation, the kind of
documents that insurers publish online for their customers. The documents were preselected and are
grouped into four areas: motor (\textit{Auto e altri veicoli}), home and family (\textit{Casa e
famiglia}), life and savings (\textit{Vita e risparmio}) and pension (\textit{Previdenza}). The
collection contains 138 \ac{pdf} documents, about 1{,}450 pages and 670{,}000 words in all. Motor
is by far the largest area, with roughly half of the pages.

The documents are of very different kinds, and each kind is written for a different purpose
(Table~\ref{tab:document-types}). The pre-contractual documents (\ac{dip}, the additional \ac{dip}
and, for investment products, the key information document) summarise a product in a fixed
structure: what is insured, what is not, the limits, the obligations. The \ac{cga} are the full
contract, organised in sections and numbered articles; they are the longest documents and the ones
that settle most questions. The pension documents describe the supplementary pension plan: costs,
investment options, taxation, advances and annuities. The forms are short, but they answer a
frequent practical question, namely what has to be sent, and how, to open a claim or ask for a
payment. Smaller groups are the privacy notices, the rules of promotional campaigns, a report on
the segregated funds and the list of the insurers designated by the road victims guarantee fund.

\begin{table}[H]
  \centering
  \small
  \renewcommand\arraystretch{1.25}
  \begin{tabularx}{\textwidth}{@{}L{4.2cm}rY@{}}
    \toprule
    \textbf{Type} & \textbf{Docs} & \textbf{What it contains} \\
    \midrule
    Forms & 49 & Claim reports, requests for payment, statements, mandates, sample letters \\
    \acs{dip} and additional \acs{dip} & 37 & Fixed-structure summary of a product \\
    \acs{cga} & 21 & The full contract, in sections and articles \\
    Pension plan documents & 15 & Costs, investment options, taxation, advances, annuities \\
    Privacy notices & 6 & How personal data is processed, for each kind of product \\
    Promotions & 5 & Rules of campaigns, discounts and cashback \\
    Key information documents & 3 & Risks, costs and scenarios of investment options \\
    Other & 2 & Report on the segregated funds, list of designated insurers \\
    \midrule
    FAQ & 170 & Short questions and answers, in five categories \\
    \bottomrule
  \end{tabularx}
  \caption{The public documents of the knowledge base.}
  \label{tab:document-types}
\end{table}

The second public source is a list of frequently asked questions, collected in a structured
file: 170 questions, each with its answer, the address of its page and one or more of
five categories. The categories reflect what customers ask most: vehicle insurance (68 questions)
and the management of a policy (60) come first, followed by claims (30), life, investment and
pension products (24) and home, family, pets and travel (18). The FAQs are written in the words
customers use, which makes them valuable for the search even when the same information is also
somewhere in a contract.

The third source is of a different nature: the data of an individual customer, such as their
policies, limits, deductibles, claims and instalments. It never enters the knowledge base. It is
read through the policy server, only after login and only for the customer who is logged in
(Section~\ref{sec:components}).

\section{Synthetic Data}
\label{sec:synthetic-data}\index{synthetic data}

All the documents used in this project are synthetic. They were not invented from scratch: they
were generated with \ac{ai} tools starting from real public insurance documentation, and every
element that could identify a real company or person was replaced. The company that results,
Velaria S.p.A., with its life insurance arm Velaria Vita and its parent group, the Gruppo Altaire,
does not exist, and neither do its partners, its products or its websites. Every document begins
with a note stating that it is synthetic and generated with \ac{ai}, and the FAQ file carries the
same note.

There are two reasons for this choice. The first is that the project, its report and its public
repository must not expose, or be traceable to, a real company: the answers of a prototype could
be mistaken for statements about that company's products. The second is that real documents, even
public ones, contain the names and personal details of real people, such as the members of a board
or the manager of a pension fund, which have no place in a student project.

What was changed is everything that points to someone: the names of the company, of the group and
of its partners, the names of other insurers, people and their dates of birth, street addresses,
phone and fax numbers, e-mail addresses and websites, tax and registration codes, bank
coordinates, and the names of products and funds that would reveal their origin. The names of
companies, products and websites were replaced with invented ones, because the assistant needs a
company to talk about; the invented websites use the \texttt{.example} domain, which is reserved for
examples and cannot belong to anyone. On the advice of the supervisor, the data that concern people
and contacts were instead replaced with placeholders in square brackets, rather than with invented
but realistic values: names and surnames become \texttt{[NOME] [COGNOME]}, places and dates of birth
\texttt{[LUOGO DI NASCITA]} and \texttt{[DATA DI NASCITA]}, and street addresses, phone and fax numbers
and e-mail addresses \texttt{[INDIRIZZO]}, \texttt{[TELEFONO]}, \texttt{[FAX]} and \texttt{[EMAIL]}.
A placeholder can never be mistaken for a real person or a real number, and it shows at a glance
where something was removed.
What was deliberately kept is the insurance content: guarantees, exclusions, limits, deductibles,
procedures and deadlines are as in the originals, because the assistant has to be tested on
realistic contracts. Public authorities such as \ac{ivass} and CONSAP, and the laws and regulations
cited by the documents, were also kept: they do not identify the company and they are part of the
answers a customer needs, for instance when explaining where to send a complaint.

Keeping the insurance content does not tie the documents to a particular company, because most of
it is common to the Italian market. The pre-contractual documents follow formats set by European
and Italian regulation, with the same sections in the same order for every insurer; the terms and
conditions follow the guidelines of the Italian insurers' association (ANIA) for simple and clear
contracts; and the procedures, such as reporting a road accident or claiming from the road victims
guarantee fund, are set by law. What made the documents recognisable were the names, the brands of
products and services and the contact details, and these are exactly what was replaced.

One limit remains. Some aggregate financial figures, such as the solvency ratio or the returns of
the funds, are still the original ones. They concern no person, but in principle they could help
someone trace the source, and they are the next candidates for being altered.

\section{From PDF to Text: Mistral OCR}
\label{sec:parsing}\index{Mistral OCR}\index{parsing}

The \acsp{pdf} of the collection are digital documents, not scans: each of them already contains its
text. That text, however, has no structure. It is a sequence of lines in the order in which they
were drawn, so titles look like ordinary sentences, the columns of a page can be mixed together,
and a table becomes a list of loose words and numbers. A search engine can still find words in it,
but it cannot cut it into meaningful passages or tell which number belongs to which row.

For this reason the documents were converted with Mistral \ac{ocr}~\cite{mistral2026ocr}, which reads
the layout of each page and returns Markdown: titles become headings, lists
become lists, tables become tables in the text, and the header and footer of each page are kept
apart from the body. The conversion was done through the Document~\ac{ai} page of Mistral's
console, on the free plan, in batches of at most ten files and five hundred pages to respect its
limits. Since the documents are public, sending them to an
external service raised no privacy issue; this is also why only public documents were sent, as on the
free plan the provider may use the inputs to improve its models. Every page of the result
begins with a marker that records its number, so that each passage can later be cited with its
page.

An \ac{ocr} model can misread a word or a number, and in insurance a wrong number is the worst
possible error. The output was therefore corrected automatically against the text layer of the
original \acsp{pdf}: where Mistral's text and the \ac{pdf}'s text disagree on a word or a number,
the \ac{pdf}, which is exact, wins. The correction script records each of its 853 changes, so that
they can be checked. Measured against the text of the \acsp{pdf}, the wrong or invented words went
from 1.16 to about 0.45 every thousand, and the numbers of the \acsp{pdf} missing from the text
from 365 to 66.

\section{Images, Tables and Other Content}
\label{sec:non-text}\index{images}\index{tables}

Not everything in a \ac{pdf} is running text, and each kind of content needed a decision.

\paragraph{Photos and images.} The knowledge base contains text only. The extraction of images was
switched off during the conversion, and the placeholders that marked their position were removed
afterwards. The images in these documents are photos, illustrations, icons and logos: they decorate
the page but carry no information that the assistant needs to answer, and the logos would even
carry the brand that the synthetic data is meant to remove. The assistant never shows images
either: when a customer needs to see the original page, the citation points to it.

\paragraph{Charts.} The text inside a chart, such as the labels and the values of a bar chart, is
recovered from the text layer of the \ac{pdf}, in the order in which it was drawn. The numbers are
there, but it is not always clear which label each one belongs to. Answers based on a chart are
therefore less reliable, and the assistant always cites the page so that the customer can check it.

\paragraph{Tables.} Tables are kept as tables, in Markdown, next to their titles, so that every
value keeps the name of its row and of its column.

\paragraph{Check boxes.} In the \acsp{pdf} the boxes of the forms and of some comparison tables are
drawings, not characters, and after the conversion they all look empty. For the forms this does not
matter, because they are blank templates; for the few tables that use ticks to show what a plan
includes, the original page remains the reference.

\paragraph{Headers, footers and notes.} The headers and footers repeat the same lines on every page
and would only add noise to the search, so they are kept apart from the body text. Footnote marks
written as superscripts were normalised, and the page markers are preserved for the citations.

\section{Manual Anonymization and Automated Substitution}
\label{sec:anonymization}\index{anonymization}

The anonymization was first carried out by hand by the team. Working on a set of base documents,
the team identified every element that points to a company or a person, decided which elements had
to be kept instead, and fixed how each reference is replaced, always in the same way for the same
entity: with a fictional name for companies and products, and with a placeholder for people and
contact details. This manual work produced the outline of the whole anonymization: the kinds of
reference to look for and the rule to apply to each of them.

The team then wrote a script that applies the same replacements to all the files of the
collection. The documents are not rewritten: rewriting a contract, by hand or with a language
model, can drop a sentence, rephrase a clause or change a limit without anyone noticing, while a
deterministic script can be checked and run again. Following the outline, the team searched the
collection for every kind of reference: names of companies and products, people, addresses, phone numbers, e-mail addresses
and websites, tax and registration codes, bank coordinates. Each one was written as a rule in a
dictionary, together with the value that replaces it and a short note; the dictionary has 209
rules. Some rules are simple words, others need care: the name of the
parent group, for instance, is also an ordinary Italian word that appears in the expression
\textit{Condizioni Generali}, so the rule replaces it only where it names the group; a phone number
is recognised whatever its punctuation; an address is recognised in its different spellings.
The script applies all the rules to the original text in a single pass. Where two rules
overlap, the longer one wins, so that a multi-word name is replaced as a whole, and a replaced name
is never replaced a second time. The text outside the substitutions does not change by a single
character, which guarantees that every limit, deductible and article number is exactly as in the
original.

The script writes a log with every substitution, its file and its page, 9{,}919 for the
documents and 747 for the FAQs. The FAQs were processed in the same way, including the
addresses of their pages; since each FAQ carries a hash of its text, the hashes were computed again
on the anonymized text. At the end a separate check searches the result for anything that is still
there: the real names of the dictionary, e-mail addresses, websites and phone numbers that are
neither placeholders nor of a public authority, tax codes, bank coordinates and dates of birth. The
check found nothing, and a comparison of headings, tables and page markers confirmed that the
structure of every document is unchanged. Because the dictionary and the log contain the real
names, they are kept in a private folder of the team and never in the public repository.

% ---------------------------------------------------------------------------------------------
% SOSPESO: deduplicazione, chunking e aggiornamento non sono ancora stati fatti.
% Il testo resta qui sotto, escluso dal documento, per quando ci arriveremo.
% ---------------------------------------------------------------------------------------------
\iffalse
\section{Deduplication}
\label{sec:deduplication}\index{duplicates}

Some documents are relevant to more than one area of the knowledge base: the form to claim a
benefit from a life policy, for instance, concerns both life and savings and pension, and the
privacy notices apply to every life product. When the collection was assembled area by area, such a
document ended up in each of the areas it belongs to. Indexing every copy would be harmful: the
search would return the same passage several times for one question, crowding out other useful
passages, and it would give that document more weight than the others.

The documents were therefore compared by their content, byte by byte, regardless of their file
names. 19 groups of identical documents were found; one copy of each group was kept and the other
23 were removed, which brought the collection from 161 to 138 documents. Among the copies, the one
kept is the one whose name says what the document is: a form called only \textit{allegato 3}, for
example, was also present as \textit{modulo unico precontrattuale}, and the second name was kept.
The areas are not lost: each document keeps the list of all the areas it belongs to, and this list
will become part of its metadata, so that a form kept under life and savings is still found when a
question concerns pension. The same check also revealed one file with a wrong title, identical to
another form, which was removed.

Documents that are similar but not identical, such as the \acp{dip} of car and motorcycle
insurance, were all kept. They share many sentences, but they describe different products, and
merging them would give wrong answers on the details that differ.

\section{Chunking and Metadata}
\label{sec:chunking}\index{chunking}\index{metadata}\index{edition of the conditions}\index{Docling}\index{OpenSearch}\index{hybrid search}

Chunking is the next step of the work; this section describes how it is planned. A chunk is the
passage that the search returns and the assistant cites. It must be small enough to be precise and
large enough to make sense on its own; the right size depends on how each type of document is
written, so the documents will not be cut at a fixed length but along their own structure
(Table~\ref{tab:document-types}). The articles of the \ac{cga} are the natural unit: an article
regulates one matter and is what a customer is referred to. The \acp{dip} and the pension documents
will be divided by section, the forms kept whole or divided by section, and a FAQ will be one chunk
made of the question together with its answer, because the question is written in the same words a
customer would use. The planned tool is the chunker of Docling~\cite{docling2026catalog}, which can
follow this structure from the Markdown headings produced by Mistral \ac{ocr} and divide a section
only when it exceeds the length that the embedding model handles well, at the boundary between
paragraphs and using the same tokenizer as the model. Each chunk will start with the titles of the
sections it belongs to, so that a passage such as \enquote{the deductible is 10\%} still says which
guarantee it refers to.

Each chunk will carry metadata that serve two purposes: filtering the search and citing the source.
They are the type of document, the product, the areas where it is published (already recorded
during the deduplication), the edition of the conditions and the dates in which it is valid, the
section and article, the page, taken from the page markers, and the source file; a FAQ also keeps
the address of its page. With them the assistant can cite a passage as \enquote{CGA Auto
ed.~03/2024, art.~12, p.~18}, and the citation opens the original document, which is kept on the
server.

The edition deserves a word of its own. Insurance conditions change over time, but a policy remains
bound to the edition that was signed. All editions will therefore stay in the index: when a logged-in
customer asks about their policy, the search is restricted to the edition of that policy, while a
general question about a product reads the edition currently on sale. Issued policies, which are
personal, never enter this shared index.

The chunks will be indexed in OpenSearch both as vectors, computed with Qwen3-Embedding, and as
text for a \ac{bm25} keyword search, and the two searches combined. The vector search finds passages
that say the same thing in other words; the keyword search finds exact terms, such as the name of a
guarantee or the number of a form, that a vector may blur.

\section{Keeping the Knowledge Base Up to Date}
\label{sec:updates}\index{ingestion}

Documents change: a new edition of the conditions comes out, a form is replaced, a FAQ is added.
Once the first index is built, a scheduled script will look for new or updated documents and pass
them through the same steps as the first collection: conversion, correction, anonymization,
deduplication and chunking. The new chunks will be written into a new index, which is evaluated on
a fixed set of test questions before it replaces the one in use; the switch is a change of alias,
so going back to the previous index takes an instant. The FAQ entries that come from the operators'
corrections, once approved by the team (Section~\ref{sec:corrections}), will enter the knowledge
base in the same way.
\fi
